% ============================================================
% 纯答案册·课时18-2.8②压轴综合二 —— 工具/答案抽册器.py 抽册生成件（勿手改；第三档＝纯答案单独成册）
% 源件（只读）：工作区/M3-第2章量产0913/成卷/导学件/课时18-2.8②压轴综合二/main.tex（md5 bee6b15f279bfb56f9fdeac770b128db）
%% 口径：题面不重印；逐题＝题号(印面号同正册)＋[答案]＋[详解]，值/详解照源件逐字
%       （钉值门硬断言：册值≡正册件内值≡值台账值tex，逐字全等）。册序＝正册装配序＝印面号 1..21 升序（同序同号）；键序断言基准＝题面库冻结 manifest canonical 键序。
%% 三档导出：整体 main-true／纯题 main-false（在产）｜本册＝纯答案册（本件）
%% 双档：ansbook-true.tex＝详解印本；ansbook-false.tex＝纯值速查本（详解不印）
% 编译：xelatex <壳> ×2，cwd＝本目录；sty＝qp-m3.sty 本地挂载（md5 7c3930362be8a0a2bdf21bbf8ac16573，锁校验）
% ============================================================
\documentclass[fontset=none]{ctexart}
\usepackage{qp-m3}
% —— 印面逐字符号路由补钉（承源件；− 走 TNR、▱ NSC 子块；禁改 sty 保 md5） ——
\xeCJKDeclareCharClass{Default}{"2212}%
\setCJKfamilyfont{nscblk}[Path=C:/提示词/工作区/字替对照-0909/variantF/fonts/,BoldFont=NSC-w600.ttf]{NSC-w500.ttf}%
\xeCJKDeclareSubCJKBlock{nscblk}{"25B1}%
\newcommand{\pxparallelogram}{{\CJKfamily{nscblk}▱}}%
% —— 断行微调（承母版实测档，三零门承重；\emergencystretch 不另设） ——
\xeCJKsetup{CJKglue={\hskip 0.10em plus 0.08em minus 0.05em}}%
% —— 纯答案册全线括线（长详解可跨栏断，承练习件全线括线口径；灰底盒不可跨栏断） ——
\ansblockgrayfalse
% —— 双档开关（false 档＝纯值速查：答案印、详解不印；件面开关，不动 sty 本体） ——
\newif\ifansbookdetail
\ifdefined\ansbookdetail
  \ifnum\ansbookdetail=0 \ansbookdetailfalse\else\ansbookdetailtrue\fi
\else
  \ansbookdetailtrue
\fi
% —— 页脚件名（承源件章名；件名＝<件型>答案册） ——
\renewcommand{\qpzhangming}{第二章\quad 平面解析几何}%
\renewcommand{\qpjianming}{导学件答案册}%

\begin{document}
\zhangtitle{第二章\quad 平面解析几何}
\jietitle{2.8 直线与圆锥曲线的位置关系（二）}
\par\glueguard{1}\addvspace{2pt}\noindent\biaoqian{【纯答案册】}{\fangsong 题面不重印\quad 印面号与正册同号同序}\par\addvspace{2pt}

\begin{multicols}{2}
\emergencystretch=1em

\begin{ansblock}[2章-练-课时18-01]
% ans:2章-练-课时18-01
\ansitem{1}{B（|MF₂|=5/2）}
\ifansbookdetail
\ansnote{详解}{\(F_1(-2,0)\)、\(F_2(2,0)\)、\(A(-1,0)\)，故 \(F_1A/F_1F_2=1/4\)．\(NA\parallel MF_2\)，得 \(\triangle F_1AN\sim\triangle F_1MF_2\)，故 \(F_1N/F_1M=1/4\)．设 \(M(x_0,y_0)\)（\(x_0>0\)，\(y_0>0\)），则 \(N=F_1+(1/4)(M-F_1)=((x_0-6)/4,y_0/4)\)．\(M\)、\(N\) 均在 \(E\) 上：\(x_0^{2}-y_0^{2}/3=1\)，且 \(((x_0-6)/4)^{2}-(y_0/4)^{2}/3=1\)，后者乘16得 \((x_0-6)^{2}-y_0^{2}/3=16\)；两式相减 \((x_0-6)^{2}-x_0^{2}=15\)，即 \(-12x_0+36=15\)，解得 \(x_0=7/4\)，进而 \(y_0^{2}=3(x_0^{2}-1)=99/16\)，\(y_0=3\sqrt{11}/4\)．于是 \(|MF_2|=\sqrt{(7/4-2)^{2}+99/16}=\sqrt{100/16}=5/2\)．故选：B．}
\fi
\end{ansblock}

\begin{ansblock}[2章-练-课时18-07]
% ans:2章-练-课时18-07
\ansitem{2}{(1) 4/5；(2) 恒过定点 (0,±10)}
\ifansbookdetail
\ansnote{详解}{(1) 由 \(2b=a+c\) 得 \(c=2b-a\)，代入 \(c^{2}=a^{2}-b^{2}\)：\((2b-a)^{2}=a^{2}-b^{2}\)，展开得 \(4b^{2}-4ab=0\)，故 \(b/a=4/5\)．\par\noindent (2) \(2c=12\)，故 \(c=6\)；由 \(b^{2}=16a^{2}/25\) 与 \(a^{2}-b^{2}=36\) 得 \(9a^{2}/25=36\)，故 \(a^{2}=100\)、\(b^{2}=64\)，\(A(0,8)\)．设 \(P(x_0,y_0)\)（\(x_0\neq 0\)），则 \(Q(-x_0,-y_0)\)．直线 \(AP\)：\(y=(y_0-8)x/x_0+8\)，令 \(y=0\) 得 \(M(8x_0/(8-y_0),0)\)；同理 \(N(-8x_0/(8+y_0),0)\)．以 \(MN\) 为直径的圆：\((x-8x_0/(8-y_0))(x+8x_0/(8+y_0))+y^{2}=0\)，其 \(x\) 的交叉项 \(=8x_0[(8-y_0)-(8+y_0)]/(64-y_0^{2})=-16x_0y_0/(64-y_0^{2})\)．由 \(x_0^{2}/100+y_0^{2}/64=1\) 得 \(64-y_0^{2}=64x_0^{2}/100=16x_0^{2}/25\)，故圆方程化为 \(x^{2}+y^{2}-(25y_0/x_0)x-100=0\)．令 \(x=0\) 得 \(y=\pm 10\)，与 \(P\) 无关，故恒过定点 \((0,10)\) 与 \((0,-10)\)（\(P\) 异于 \(A\) 保证 \(x_0\neq 0\)、\(64-y_0^{2}>0\)）．}
\fi
\end{ansblock}

\begin{ansblock}[2章-练-课时18-08]
% ans:2章-练-课时18-08
\ansitem{3}{(1) x²/4+y²/3=1；(2) 存在，R(4,0)}
\ifansbookdetail
\ansnote{详解}{(1) 圆 \(A\) 即 \((x+1)^{2}+y^{2}=16\)，圆心 \(A(-1,0)\)、半径4．由垂直平分线性质得 \(|NM|=|NB|\)，而 \(N\) 在 \(MA\) 上，故 \(|NA|+|NB|=|NA|+|NM|=|MA|=4>|AB|=2\)，轨迹为以 \(A\)、\(B\) 为焦点的椭圆：\(2a=4\)、\(2c=2\)，即 \(a^{2}=4\)、\(b^{2}=3\)，故 \(C\)：\(x^{2}/4+y^{2}/3=1\)．\par\noindent (2) 设 \(R(t,0)\)．联立得 \((4k^{2}+3)x^{2}-8k^{2}x+4k^{2}-12=0\)，故 \(x_1+x_2=8k^{2}/(4k^{2}+3)\)、\(x_1x_2=(4k^{2}-12)/(4k^{2}+3)\)．\(\angle ORP=\angle ORQ\) 等价于 \(k_{RP}+k_{RQ}=0\)，即 \(y_1/(x_1-t)+y_2/(x_2-t)=0\)．以 \(y_i=k(x_i-1)\) 代入，分子 \(=k[2x_1x_2-(1+t)(x_1+x_2)+2t]=0\)；\(k\neq 0\) 时即 \(2(4k^{2}-12)-8k^{2}(1+t)+2t(4k^{2}+3)=0\)，展开得 \(6t-24=0\)，故 \(t=4\)（\(k=0\) 时 \(P\)、\(Q\) 为长轴两端点，等式显然成立）．故存在 \(R(4,0)\)，与 \(k\) 无关．}
\fi
\end{ansblock}

\begin{ansblock}[2章-练-课时18-10]
% ans:2章-练-课时18-10
\ansitem{4}{B}
\ifansbookdetail
\ansnote{详解}{过焦点的弦分两型：跨支型（\(A\)、\(B\) 各在一支）弦长下限为 \(2a\)，且实轴所在直线（长 \(2a\)）恰1条，长 \(L>2a\) 时对称地有2条；单支型（同在右支）弦长下限为通径 \(2b^{2}/a=6/a\)，通径（\(l\perp x\) 轴）恰1条，\(L>6/a\) 时有2条．且 \(6/a\) 与 \(2a\) 不可能同时等于6（\(a=1\) 时 \(6/a=6\)、\(2a=2\)；\(a=3\) 时 \(2a=6\)、\(6/a=2\)）．令 \(L=6\) 分段计数：\(a<1\) 时跨支 \(6>2a\) 有2条、单支 \(6<6/a\) 有0条，合计2，符合；\(a=1\) 时 \(2+1=3\)，不合；\(1<a<3\) 时 \(2+2=4\)，不合；\(a=3\) 时 \(1+2=3\)，不合；\(a>3\) 时跨支 \(6<2a\) 有0条、单支 \(6>6/a\) 有2条，合计2，符合．故 \(a\in(0,1)\cup(3,+\infty)\)，选 B．}
\fi
\end{ansblock}

\begin{ansblock}[2章-练-课时18-11]
% ans:2章-练-课时18-11
\ansitem{5}{(1) x²/4+y²/2=1；(2) 不存在（QA·QB≡−15/16 恒<1）}
\ifansbookdetail
\ansnote{详解}{(1) \(e=\sqrt{2}/2\)，故 \(a^{2}=2c^{2}\)、\(b^{2}=c^{2}\)；三角形面积 \(=(a-c)b/2=(\sqrt{2}-1)c^{2}/2=\sqrt{2}-1\)，得 \(c^{2}=2\)，故 \(a^{2}=4\)、\(b^{2}=2\)，椭圆 \(C\)：\(x^{2}/4+y^{2}/2=1\)．\par\noindent (2) 直线过 \((1,0)\)（椭圆内一点），恒有两交点．联立得 \((1+2k^{2})x^{2}-4k^{2}x+2k^{2}-4=0\)，故 \(x_1+x_2=4k^{2}/(1+2k^{2})\)、\(x_1x_2=(2k^{2}-4)/(1+2k^{2})\)；\(y_1y_2=k^{2}(x_1-1)(x_2-1)=-3k^{2}/(1+2k^{2})\)．于是 \(QA\cdot QB=(x_1-7/4)(x_2-7/4)+y_1y_2=x_1x_2-(7/4)(x_1+x_2)+49/16+y_1y_2\)，代入化简 \(=[(2k^{2}-4)-7k^{2}-3k^{2}]/(1+2k^{2})+49/16=(-8k^{2}-4)/(1+2k^{2})+49/16=-4+49/16=-15/16\)，与 \(k\) 无关且恒小于1，故不存在这样的 \(k\)．}
\fi
\end{ansblock}

\begin{ansblock}[2章-练-课时18-09]
% ans:2章-练-课时18-09
\ansitem{6}{(1) 两条件均得 x²/4+y²/3=1；(2) DE 恒过定点 (8/5,0)}
\ifansbookdetail
\ansnote{详解}{(1) 条件①：\(e=1/2\)，故 \(a=2c\)；焦点到相应准线距离 \(a^{2}/c-c=3\)，即 \(4c-c=3\)，得 \(c=1\)，故 \(a^{2}=4\)、\(b^{2}=3\)．条件②：\(\odot M\) 圆心 \((6,0)\)、半径4，椭圆在圆内一侧外切，切点必为右顶点，故 \(a+4=6\)，\(a=2\)；\(\odot N\) 圆心 \((0,2\sqrt{3})\)、半径 \(\sqrt{3}\)，故 \(b=\sqrt{3}\)、\(b^{2}=3\)．回验唯一性（外切即恰一公共点）：\(\odot M\) 与椭圆联立消 \(y\) 得 \(x^{2}-48x+92=0\)，\(x=2\) 或 \(x=46\)（46越界，舍）；\(\odot N\) 联立得 \(y^{2}+12\sqrt{3}y-39=0\)，\(y=\sqrt{3}\) 或 \(y=-13\sqrt{3}\)（越界，舍）．两条件同得 \(x^{2}/4+y^{2}/3=1\)．\par\noindent (2) \(F(1,0)\)．设 \(A(x_1,y_1)\)（\(y_1>0\)）、\(B(-x_1,-y_1)\)，满足 \(3x_1^{2}+4y_1^{2}=12\)．直线 \(AF\)：\(x=my+1\)（\(m=(x_1-1)/y_1\)），直线 \(BF\)：\(x=ny+1\)（\(n=(x_1+1)/y_1\)），则 \(n-m=2/y_1\)．\(x=my+1\) 代入 \(3x^{2}+4y^{2}=12\)：\((3m^{2}+4)y^{2}+6my-9=0\)，故 \(y_1y_D=-9/(3m^{2}+4)\)，即 \(y_D=-9/[(3m^{2}+4)y_1]\)；同理 \(y_E=9/[(3n^{2}+4)y_1]\)，且 \(x_D=my_D+1\)、\(x_E=ny_E+1\)．\par\noindent \(D\)、\(E\)、\(P(8/5,0)\) 三点共线等价于 \(y_D(x_E-8/5)=y_E(x_D-8/5)\)，即 \(y_Dy_E(n-m)=(3/5)(y_D-y_E)\)（记为①）．代入 \(y_Dy_E=-81/[(3m^{2}+4)(3n^{2}+4)y_1^{2}]\) 与 \(y_D-y_E=-(9/y_1)(3m^{2}+3n^{2}+8)/[(3m^{2}+4)(3n^{2}+4)]\)，则①等价于 \(30/y_1^{2}=3(m^{2}+n^{2})+8\)，即 \(30=3y_1^{2}(m^{2}+n^{2})+8y_1^{2}=3[(x_1-1)^{2}+(x_1+1)^{2}]+8y_1^{2}=6x_1^{2}+6+8y_1^{2}=2(3x_1^{2}+4y_1^{2})+6=24+6=30\)，恒成立．故 \(DE\) 恒过定点 \((8/5,0)\)．}
\fi
\end{ansblock}

\begin{ansblock}[2章-练-课时18-03]
% ans:2章-练-课时18-03
{\catcode`\_=12 \ansitem{7}{(1) y²=4x、x²/4+y²/3=1；(2) S_max=2√3/3}}
\ifansbookdetail
\ansnote{详解}{(1) 焦点到准线的距离 \(p=2\)，故抛物线 \(y^{2}=4x\)，\(F(1,0)\)；椭圆 \(c=1\)、\(e=c/a=1/2\)，故 \(a=2\)、\(b^{2}=a^{2}-c^{2}=3\)，椭圆 \(x^{2}/4+y^{2}/3=1\)．\par\noindent (2) 设 \(l\)：\(x=my+n\)（不过 \(F\)，故 \(n\neq 1\)）．交抛物线：\(y^{2}-4my-4n=0\)，得 \(y_1y_2=-4n\)、\(x_1x_2=(y_1y_2)^{2}/16=n^{2}\)．由 \(OA\cdot OB=n^{2}-4n=-3\) 得 \(n=1\)（此时 \(l\) 过 \(F(1,0)\)，舍）或 \(n=3\)．交椭圆：\(3(my+3)^{2}+4y^{2}=12\)，即 \((3m^{2}+4)y^{2}+18my+15=0\)，\(\Delta=144m^{2}-240\geqslant 0\)，故 \(m^{2}\geqslant 5/3\)．\(l\) 交 \(x\) 轴于 \(T(3,0)\)，\(|TF|=2\)，故 \(S_{\triangle CDF}=|TF|\cdot|y_C-y_D|/2=|y_C-y_D|\)，\(S^{2}=(144m^{2}-240)/(3m^{2}+4)^{2}\)．令 \(u=3m^{2}+4\)（\(u\geqslant 9\)）：\(S^{2}=48/u-432/u^{2}\)；再令 \(v=1/u\in(0,1/9]\)：\(S^{2}=48v-432v^{2}\) 在 \(v=1/18\) 处取最大（\(1/18\in(0,1/9]\)），\(S^{2}\) 最大值为 \(48/18-432/324=4/3\)，故 \(S\) 的最大值为 \(2\sqrt{3}/3\)，此时 \(3m^{2}+4=18\)，\(m^{2}=14/3\)，满足 \(m^{2}\geqslant 5/3\)．}
\fi
\end{ansblock}

\begin{ansblock}[2章-练-课时18-04]
% ans:2章-练-课时18-04
\ansitem{8}{(1) d=|x₁y₂−x₂y₁|/√(x₁²+y₁²)，证毕；①S=√2；②m=−1/2（此时 S≡√2）}
\ifansbookdetail
\ansnote{详解}{(1) \(l_1\)：\(y_1x-x_1y=0\)（\(x_1=0\) 时即 \(x=0\)，公式仍适用），故 \(d=|y_1x_2-x_1y_2|/\sqrt{x_1^{2}+y_1^{2}}\)．\(O\) 同时平分 \(AB\) 与 \(CD\)，故 \(S=4S_{\triangle OAC}=4\cdot|OA|\cdot d/2=2\sqrt{x_1^{2}+y_1^{2}}\cdot|x_1y_2-x_2y_1|/\sqrt{x_1^{2}+y_1^{2}}=2|x_1y_2-x_2y_1|\)，得证．\par\noindent (2) 记 \(y_i=k_ix_i\)，则 \(x_i^{2}=1/(1+2k_i^{2})\)，\(S=2|x_1||x_2||k_2-k_1|=2|k_2-k_1|/\sqrt{(1+2k_1^{2})(1+2k_2^{2})}\)．\par\noindent ① \(k_1k_2=-1/2\)：分母 \(=(1+2k_1^{2})(1+2k_2^{2})=1+2(k_1^{2}+k_2^{2})+4k_1^{2}k_2^{2}=2(1+k_1^{2}+k_2^{2})\)；\((k_2-k_1)^{2}=k_1^{2}+k_2^{2}-2k_1k_2=1+k_1^{2}+k_2^{2}\)，故 \(S^{2}=4(1+k_1^{2}+k_2^{2})/[2(1+k_1^{2}+k_2^{2})]=2\)，即 \(S=\sqrt{2}\)．\par\noindent ② 记 \(T=k_1^{2}+k_2^{2}\)，则 \((k_2-k_1)^{2}=T-2m\)、分母 \(=2T+1+4m^{2}\)，\(S^{2}=4(T-2m)/(2T+1+4m^{2})\)．要与 \(T\) 无关，须存在 \(\lambda\) 使 \(T-2m\equiv\lambda(2T+1+4m^{2})\)：比较系数得 \(2\lambda=1\)、\(\lambda(1+4m^{2})=-2m\)，故 \((1+4m^{2})/2=-2m\)，即 \(4m^{2}+4m+1=0\)，解得 \(m=-1/2\)（与①自洽，此时对一切 \(T\) 有 \(S^{2}=2\)）．}
\fi
\end{ansblock}

\begin{ansblock}[2章-练-课时18-05]
% ans:2章-练-课时18-05
\ansitem{9}{(1) x²/4+y²/3=1；(2) 存在，t=−2/3，定值 1/2}
\ifansbookdetail
\ansnote{详解}{(1) \(y^{2}=4x\) 的焦点为 \((1,0)\)，故 \(c=1\)；长轴 \(2a=4\)，故 \(a=2\)、\(b^{2}=3\)，椭圆 \(E\)：\(x^{2}/4+y^{2}/3=1\)．\par\noindent (2) \(l\)：\(y=k(x-1)\)（\(k\neq 0\)）．交椭圆：\((3+4k^{2})x^{2}-8k^{2}x+4k^{2}-12=0\)，故 \(|AB|=\sqrt{1+k^{2}}\cdot\sqrt{(x_1+x_2)^{2}-4x_1x_2}=12(1+k^{2})/(3+4k^{2})\)．交抛物线：\(k^{2}x^{2}-(2k^{2}+4)x+k^{2}=0\)，故 \(x_3+x_4=(2k^{2}+4)/k^{2}\)，\(|MN|=x_3+x_4+2=4(1+k^{2})/k^{2}\)（焦点弦）．于是 \(2/|AB|+t/|MN|=(3+4k^{2})/(6(1+k^{2}))+tk^{2}/(4(1+k^{2}))=[(8+3t)k^{2}+6]/(12(1+k^{2}))\)，与 \(k\) 无关等价于 \(8+3t=6\)，故 \(t=-2/3\)，此时定值为 \(6/12=1/2\)．}
\fi
\end{ansblock}

\begin{ansblock}[2章-练-课时18-06]
% ans:2章-练-课时18-06
\ansitem{10}{(1) y=1 或 x=0 或 y=−x+1；(2) 仅 k₁+k₂=−4 为定值，其余三个均不是}
\ifansbookdetail
\ansnote{详解}{(1) 分三类：与对称轴（\(x\) 轴）平行的直线 \(y=1\)，恰一公共点；过顶点的切线 \(x=0\)；斜线 \(y=kx+1\)（\(k\neq 0\)）代入得 \(k^{2}x^{2}+(2k+4)x+1=0\)，令 \(\Delta=(2k+4)^{2}-4k^{2}=16k+16=0\)，得 \(k=-1\)，即 \(y=-x+1\)．共三条．\par\noindent (2) 有两交点，故 \(k\) 存在、\(k\neq 0\) 且 \(\Delta=16k+16>0\)，即 \(k>-1\)．联立 \(k^{2}x^{2}+(2k+4)x+1=0\)，得 \(x_1+x_2=-(2k+4)/k^{2}\)、\(x_1x_2=1/k^{2}\)．\(k_i=y_i/x_i=k+1/x_i\)：\(k_1+k_2=2k+(x_1+x_2)/(x_1x_2)=2k-(2k+4)=-4\)，为定值；\(k_1k_2=y_1y_2/(x_1x_2)=-4k\)，随 \(k\) 变；\(|k_1-k_2|=|x_2-x_1|/|x_1x_2|=4\sqrt{k+1}\)，随 \(k\) 变；\(k_1/k_2\) 同理随 \(k\) 变．故仅 \(k_1+k_2=-4\) 为定值．}
\fi
\end{ansblock}

\begin{ansblock}[2章-练-课时18-15]
% ans:2章-练-课时18-15
\ansitem{11}{(1) x²/4+y²=1；(2) 存在，min=2，此时 x−√2y−√3=0 或 x+√2y−√3=0}
\ifansbookdetail
\ansnote{详解}{(1) 周长 \(=|AB|+|AF_2|+|BF_2|=4a=8\)，故 \(a=2\)；\(e=\sqrt{3}/2\)，故 \(c=\sqrt{3}\)、\(b^{2}=1\)，椭圆 \(E\)：\(x^{2}/4+y^{2}=1\)．\par\noindent (2) 由相切得 \(d(O,l_2)=|m|/\sqrt{1+k^{2}}=1\)，即 \(m^{2}=1+k^{2}\)．焦半径公式（\(|x_i|\leqslant 2\) 保证非负）：\(|MF_2|=a-ex_M=2-(\sqrt{3}/2)x_M\)，同理 \(|NF_2|=2-(\sqrt{3}/2)x_N\)，故和 \(=4-(\sqrt{3}/2)(x_M+x_N)\)．联立 \(y=kx+m\) 与 \(x^{2}+4y^{2}=4\)：\((1+4k^{2})x^{2}+8kmx+4m^{2}-4=0\)，故 \(x_M+x_N=-8km/(1+4k^{2})\)；\(km<0\)，故 \(x_M+x_N=8\sqrt{k^{2}(1+k^{2})}/(1+4k^{2})\)．设 \(u=k^{2}>0\)，则 \(x_M+x_N=8\sqrt{u(1+u)}/(1+4u)\)；对 \(u(1+u)/(1+4u)^{2}\) 求导，分子正比于 \((1+2u)(1+4u)-8(u+u^{2})=1-2u\)，故 \(u=1/2\) 时最大，\(x_M+x_N\) 的最大值为 \(8\cdot(\sqrt{3}/2)/3=4\sqrt{3}/3\)，和的最小值为 \(4-(\sqrt{3}/2)\cdot(4\sqrt{3}/3)=2\)．此时 \(k^{2}=1/2\)、\(m^{2}=3/2\) 且 \(km<0\)，故 \((k,m)=(\sqrt{2}/2,-\sqrt{6}/2)\) 或 \((-\sqrt{2}/2,\sqrt{6}/2)\)，即 \(l_2\)：\(x-\sqrt{2}y-\sqrt{3}=0\) 或 \(x+\sqrt{2}y-\sqrt{3}=0\)．}
\fi
\end{ansblock}

\begin{ansblock}[2章-练-课时18-02]
% ans:2章-练-课时18-02
\ansitem{12}{CD}
\ifansbookdetail
\ansnote{详解}{\(F(2,0)\)，准线 \(x=-2\)，故 \(M(-2,0)\)，\(l\)：\(y=k(x+2)\)．联立得 \(k^{2}x^{2}+(4k^{2}-8)x+4k^{2}=0\)；由两交点得 \(k^{2}\neq 0\) 且 \(\Delta=64(1-k^{2})>0\)，即 \(-1<k<1\) 且 \(k\neq 0\)，A 漏 \(k\neq 0\)，错．\par\noindent 改以 \(y\) 为主元：\(x=y^{2}/8\) 代入 \(l\) 得 \(ky^{2}-8y+16k=0\)，故 \(y_1+y_2=8/k\)、\(y_1y_2=16\)；而 \(x_1x_2=4\)，故 \(8x_1x_2=32\neq 16\)，B 错．\par\noindent \(FA\cdot FB=(x_1-2)(x_2-2)+y_1y_2=x_1x_2-2(x_1+x_2)+4+y_1y_2\)，其中 \(x_1+x_2=(8-4k^{2})/k^{2}\)，代入得 \(FA\cdot FB=32-16/k^{2}\)，当 \(k^{2}=1/2\) 时为0，故 \(\angle AFB\) 可为直角，C 对．\par\noindent \(S=|MF|\cdot|y_1-y_2|/2=2\sqrt{(y_1+y_2)^{2}-4y_1y_2}=2\sqrt{64/k^{2}-64}\)；\(k^{2}=1/2\) 时 \(|y_1-y_2|=8\)，故 \(S=16\)，D 对．故选：CD．}
\fi
\end{ansblock}

\begin{ansblock}[2章-练-课时18-12]
% ans:2章-练-课时18-12
\ansitem{13}{(1) y²=8x；(2) 存在，t=±√2/2}
\ifansbookdetail
\ansnote{详解}{(1) 设 \(y^{2}=ax\)，准线 \(x=-a/4=-2\)，得 \(a=8\)，故抛物线 \(y^{2}=8x\)．\par\noindent (2) 联立 \(x=ty+1\)：\(y^{2}-8ty-8=0\)，得 \(y_1+y_2=8t\)、\(y_1y_2=-8\)；\(x_1+x_2=t(y_1+y_2)+2=8t^{2}+2\)、\(x_1x_2=t^{2}y_1y_2+t(y_1+y_2)+1=1\)．设 \(C(x_0,0)\)（\(x_0<0\)）．\(C\) 在以 \(AB\) 为直径的圆上，故 \(\angle ACB=90^{\circ}\)，即 \(CA\cdot CB=0\)：\((x_1-x_0)(x_2-x_0)+y_1y_2=0\)，得 \(x_0^{2}-(8t^{2}+2)x_0+(x_1x_2+y_1y_2)=0\)，即 \(x_0^{2}-(8t^{2}+2)x_0-7=0\)（记为①）．内心在 \(x\) 轴上等价于 \(x\) 轴平分 \(\angle ACB\)，即 \(k_{CA}+k_{CB}=0\)：\(y_1(x_2-x_0)+y_2(x_1-x_0)=0\)，得 \(2ty_1y_2+(y_1+y_2)(1-x_0)=0\)，即 \(-16t+8t(1-x_0)=0\)，故 \(8t(-1-x_0)=0\)．\(t=0\) 时直线 \(x=1\) 垂直于 \(x\) 轴，与题设矛盾，舍；故 \(x_0=-1<0\)，代入①：\(1+(8t^{2}+2)-7=0\)，故 \(t^{2}=1/2\)，\(t=\pm\sqrt{2}/2\)．}
\fi
\end{ansblock}

\begin{ansblock}[2章-练-课时18-13]
% ans:2章-练-课时18-13
\ansitem{14}{(1) x²/2+y²=1；(2) y=(2−√2)/2·(x+2) 或 y=0}
\ifansbookdetail
\ansnote{详解}{(1) 焦距为2，故 \(c=1\)、\(a^{2}=b^{2}+1\)；点 \((1,\sqrt{2}/2)\) 代入：\(1/a^{2}+(1/2)/b^{2}=1\)．令 \(u=b^{2}\)：\(1/(u+1)+1/(2u)=1\)，即 \(2u+(u+1)=2u(u+1)\)，整理得 \(2u^{2}-u-1=0\)，解得 \(u=1\)，故 \(a^{2}=2\)、\(b^{2}=1\)，椭圆 \(C\)：\(x^{2}/2+y^{2}=1\)．\par\noindent (2) 先取 \(y=0\)：\(A\)、\(B\) 为 \((\pm\sqrt{2},0)\)，\(G\) 在 \(y\) 轴上，\(|GA|=|GB|\) 成立，是一解．\(k\neq 0\) 时设 \(l\)：\(y=k(x+2)\)，联立 \(x^{2}+2y^{2}=2\)：\((1+2k^{2})x^{2}+8k^{2}x+8k^{2}-2=0\)，弦中点 \(M\)：\(x_M=-4k^{2}/(1+2k^{2})\)、\(y_M=k(x_M+2)=2k/(1+2k^{2})\)．\(|GA|=|GB|\) 等价于 \(GM\perp l\)，即 \(x_M+k(y_M+1/2)=0\)，两边乘 \(2(1+2k^{2})\)：\(-8k^{2}+4k^{2}+k(1+2k^{2})=0\)，即 \(k(2k^{2}-4k+1)=0\)（\(k\neq 0\)），故 \(k=(2\pm\sqrt{2})/2\)．两交点条件：\(\Delta=8-16k^{2}>0\)，即 \(k^{2}<1/2\)；\(k=(2+\sqrt{2})/2\) 时 \(k^{2}=(3+2\sqrt{2})/2>1/2\)，舍；\(k=(2-\sqrt{2})/2\) 时 \(k^{2}=(3-2\sqrt{2})/2<1/2\)，符合．故 \(l\)：\(y=(2-\sqrt{2})/2\cdot(x+2)\) 或 \(y=0\)．}
\fi
\end{ansblock}

\begin{ansblock}[2章-练-课时18-14]
% ans:2章-练-课时18-14
\ansitem{15}{(1) a=2、b=1；(2) 存在，8x+3y−8=0}
\ifansbookdetail
\ansnote{详解}{(1) \(C_2\) 与 \(x\) 轴交于 \((\pm 1,0)\)，即 \(A(-1,0)\)、\(B(1,0)\)，且 \(x^{2}/b^{2}=1\) 在 \(y=0\) 处成立，故 \(b=1\)．\(e=c/a=\sqrt{3}/2\)，故 \(c^{2}=3a^{2}/4\)；\(a^{2}-c^{2}=b^{2}=1\)，得 \(a^{2}/4=1\)，故 \(a=2\)．\par\noindent (2) \(C_1\)：\(y^{2}/4+x^{2}=1\)（\(y\geqslant 0\)）．设 \(l\)：\(y=k(x-1)\)（\(k\neq 0\)；\(k=0\) 时 \(l\) 与 \(x\) 轴重合，不合）．交 \(C_1\)：\((k^{2}+4)x^{2}-2k^{2}x+k^{2}-4=0\)，一根 \(x=1\)（点 \(B\)），故 \(x_P=(k^{2}-4)/(k^{2}+4)\)、\(y_P=k(x_P-1)=-8k/(k^{2}+4)\)．交 \(C_2\)：\(k(x-1)=-(x-1)(x+1)\)，故 \((x-1)(k+x+1)=0\)，\(x_Q=-k-1\)（\(x_Q\neq 1\)，故 \(k\neq -2\)）、\(y_Q=k(x_Q-1)=-k^{2}-2k\)．\par\noindent \(\overrightarrow{AP}=(2k^{2}/(k^{2}+4),-8k/(k^{2}+4))=(2k/(k^{2}+4))\cdot(k,-4)\)；\(\overrightarrow{AQ}=(-k,-k^{2}-2k)=-k(1,k+2)\)．以 \(PQ\) 为直径的圆过 \(A\) 等价于 \(AP\perp AQ\)，即 \(\overrightarrow{AP}\cdot\overrightarrow{AQ}=0\)：\((2k/(k^{2}+4))\cdot(-k)\cdot[k-4(k+2)]=0\)，\(k\neq 0\)，故 \(3k+8=0\)，\(k=-8/3\)．回验落位：\(y_P=48/25>0\)（\(P\) 在上半椭圆）、\(y_Q=-16/9\leqslant 0\)（\(Q\) 在下半抛物线），符合．故存在 \(l\)：\(y=-(8/3)(x-1)\)，即 \(8x+3y-8=0\)．}
\fi
\end{ansblock}

\begin{ansblock}[2章-练-课时18-16]
% ans:2章-练-课时18-16
\ansitem{16}{(1) e=2；(2) 恒过定点 (4,0) 与 (−2,0)（两点同时在圆上）}
\ifansbookdetail
\ansnote{详解}{(1) \(l\perp x\) 轴时 \(M\)、\(N\) 的坐标为 \((c,\pm b^{2}/a)\)；\(A\) 在 \(x\) 轴上，\(|AM|=|AN|\) 自动成立，故直角顶点只能是 \(A\)（若直角在 \(M\)，\(MA\) 需水平，即 \(y_M=0\)，不可能），故 \(\angle MAN=90^{\circ}\)，\(k_{AM}=\pm 1\)，即 \(b^{2}/a=a+c\)，故 \(c^{2}-a^{2}=a^{2}+ac\)，即 \(e^{2}-e-2=0\)，得 \(e=2\)（负根舍）．\par\noindent (2) \(2a=4\)，故 \(a=2\)、\(c=2a=4\)、\(b^{2}=12\)，\(C\)：\(x^{2}/4-y^{2}/12=1\)，\(A(-2,0)\)，\(x=a/2=1\)．设 \(l\)：\(x=ty+4\)（\(t\) 为实数，\(t=0\) 即 \(l\perp x\) 轴亦含；\(|t|<1/\sqrt{3}\) 时与右支交两点）．代入得 \((3t^{2}-1)y^{2}+24ty+36=0\)，故 \(y_1y_2=36/(3t^{2}-1)\)、\(y_1+y_2=-24t/(3t^{2}-1)\)、\(x_1+x_2=t(y_1+y_2)+8=-8/(3t^{2}-1)\)、\(x_1x_2=t^{2}y_1y_2+4t(y_1+y_2)+16=(-12t^{2}-16)/(3t^{2}-1)\)．\par\noindent 直线 \(AM\) 过 \(A(-2,0)\)、\(M(x_1,y_1)\)，在 \(x=1\) 处 \(y=3y_1/(x_1+2)\)，故 \(P(1,p)\)、\(Q(1,q)\)，其中 \(p=3y_1/(x_1+2)\)、\(q=3y_2/(x_2+2)\)．而 \((x_1+2)(x_2+2)=x_1x_2+2(x_1+x_2)+4=(-12t^{2}-16-16+12t^{2}-4)/(3t^{2}-1)=-36/(3t^{2}-1)\)，故 \(pq=9y_1y_2/[(x_1+2)(x_2+2)]=9\cdot[36/(3t^{2}-1)]\cdot[(3t^{2}-1)/(-36)]=-9\)，与 \(t\) 无关．\par\noindent 以 \(PQ\) 为直径的圆：\((x-1)^{2}+(y-p)(y-q)=0\)，即 \((x-1)^{2}+y^{2}-(p+q)y+pq=0\)．令 \(y=0\)：\((x-1)^{2}=-pq=9\)，得 \(x=4\) 或 \(x=-2\)，即对每一条满足题设的 \(l\)，\((4,0)\) 与 \((-2,0)\) 同时在该圆上．特例回验：\(t=0\) 时 \(M\)、\(N\) 为 \((4,\pm 6)\)，\(P(1,3)\)、\(Q(1,-3)\)，圆 \((x-1)^{2}+y^{2}=9\)，两点均在圆上；\(3t^{2}-1=0\) 时 \(l\) 平行于渐近线仅一交点，不在题设内．}
\fi
\end{ansblock}

\begin{ansblock}[2章-导-课时18-G1]
% ans:2章-导-课时18-G1
\ansitem{17}{C}
\ifansbookdetail
\ansnote{详解}{\(F(p/2,0)\)，准线 \(x=-p/2\)，记准线与 \(x\) 轴交点 \(A(-p/2,0)\)，\(|AF|=p\)．\(l\) 过 \(F\) 与准线交于 \(P\)，弦 \(MN\) 过 \(F\)，沿直线点序为 \(P\)、\(N\)、\(F\)、\(M\)（\(N\) 近 \(F\)、\(M\) 远 \(F\)，\(|PN|<|PF|<|PM|\)）．由 \(|PM|=2|PF|\) 及 \(|PM|=|PF|+|FM|\) 得 \(|FM|=|PF|\)，即 \(F\) 为 \(PM\) 中点．过 \(M\) 作 \(MB\perp\) 准线于 \(B\)，由定义 \(|FM|=|BM|\)；\(A\) 为 \(PB\) 中点（\(P\)、\(B\) 同在准线上，纵坐标和为0）、\(F\) 为 \(PM\) 中点，故 \(AF\) 为 \(\triangle PMB\) 中位线，\(|BM|=2|AF|=2p\)，故 \(|FM|=|PF|=2p\)．\(P\) 与 \(F\) 的水平距离为 \(p\)，设倾斜角 \(\theta\)，则 \(|PF|\cdot|\cos\theta|=p\)，故 \(2p|\cos\theta|=p\)，\(|\cos\theta|=1/2\)，\(\theta\in(0,\pi)\)，得 \(\theta=\pi/3\) 或 \(\theta=2\pi/3\)．故选：C．\par\noindent 特例回验：取 \(p=2\)（\(y^{2}=4x\)）、\(\theta=\pi/3\)，\(l\)：\(y=\sqrt{3}(x-1)\)，联立得 \(3x^{2}-10x+3=0\)，\(M(3,2\sqrt{3})\)、\(N(1/3,-2\sqrt{3}/3)\)、\(P(-1,-2\sqrt{3})\)；\(|PF|=\sqrt{4+12}=4\)、\(|PM|=\sqrt{16+48}=8\)，\(|PM|=2|PF|\) 且 \(|FM|=x_M+1=4=|PF|\)，符合；\(\theta=2\pi/3\) 与 \(\pi/3\) 关于 \(x\) 轴对称，同成立．}
\fi
\end{ansblock}

\begin{ansblock}[2章-导-课时18-G2]
% ans:2章-导-课时18-G2
\ansitem{18}{D}
\ifansbookdetail
\ansnote{详解}{\(a=2\)、\(b=\sqrt{3}\)，故 \(c=\sqrt{4-3}=1\)，\(F_2(1,0)\)．通径：\(x=1\) 代入得 \(y^{2}=3(1-1/4)=9/4\)，\(|y|=3/2\)，故 \(|AB|=3\)（\(=2b^{2}/a=3\)，两法互证）．以 \(|F_1F_2|=2\) 为底：\(S_{\triangle ABF_1}=|AB|\cdot 2/2=3\)．周长：\(|AF_1|+|BF_1|=(2a-|AF_2|)+(2a-|BF_2|)=8-|AB|=5\)，故周长为 \(5+3=8\)．内切圆半径 \(r=2\times 3/8=3/4\)．故选：D．}
\fi
\end{ansblock}

\begin{ansblock}[2章-导-课时18-G3]
% ans:2章-导-课时18-G3
\ansitem{19}{C}
\ifansbookdetail
\ansnote{详解}{设直线 \(y=kx-2\)（过 \((0,-2)\)），\(A(x_1,y_1)\)、\(B(x_2,y_2)\)．代入 \(y^{2}=8x\)：\(k^{2}x^{2}-4(k+2)x+4=0\)，\(\Delta=16(k+2)^{2}-16k^{2}=64(k+1)>0\)，故 \(k>-1\)．中点横坐标 \((x_1+x_2)/2=2(k+2)/k^{2}=2\)，故 \(k^{2}=k+2\)，即 \(k^{2}-k-2=0\)，得 \(k=2\) 或 \(k=-1\)；\(k=-1\) 时 \(\Delta=0\)（相切不成弦），舍，故 \(k=2\)．此时 \(x^{2}-4x+1=0\)，\(x_1+x_2=4\)、\(x_1x_2=1\)，弦长 \(|AB|=\sqrt{1+k^{2}}\cdot\sqrt{(x_1+x_2)^{2}-4x_1x_2}=\sqrt{5}\cdot\sqrt{12}=2\sqrt{15}\)．故选：C．}
\fi
\end{ansblock}

\begin{ansblock}[2章-导-课时18-G4]
% ans:2章-导-课时18-G4
\ansitem{20}{81π/16}
\ifansbookdetail
\ansnote{详解}{\(a=4\)、\(b=3\)，故 \(c=\sqrt{16-9}=\sqrt{7}\)，\(F(\sqrt{7},0)\)．\(x=\sqrt{7}\) 代入：\(y^{2}=9(1-7/16)=81/16\)，\(|y|=9/4\)，故 \(|AB|=9/2\)（与通径 \(2b^{2}/a=18/4=9/2\) 一致）．圆半径 \(r=|AB|/2=9/4\)，面积 \(S=\pi r^{2}=81\pi/16\)．}
\fi
\end{ansblock}

\begin{ansblock}[2章-导-课时18-G5]
% ans:2章-导-课时18-G5
\ansitem{21}{9/4}
\ifansbookdetail
\ansnote{详解}{\(2p=3\)，故 \(p=3/2\)，\(F(3/4,0)\)，\(k=\tan 30^{\circ}=\sqrt{3}/3\)，直线 \(y=(\sqrt{3}/3)(x-3/4)\)．代入 \(y^{2}=3x\)：\((x-3/4)^{2}=9x\)，即 \(x^{2}-(21/2)x+9/16=0\)（\(\Delta=441/4-9/4=108>0\)），\(x_1+x_2=21/2\)．由抛物线定义 \(|AB|=x_1+x_2+p=21/2+3/2=12\)．\(O\) 到直线 \(AB\)（即 \(\sqrt{3}x-3y-3\sqrt{3}/4=0\)）的距离 \(d=(3\sqrt{3}/4)/\sqrt{3+9}=3/8\)，故 \(S_{\triangle ABO}=|AB|\cdot d/2=12\times(3/8)/2=9/4\)．}
\fi
\end{ansblock}

% ---- 尾块（\tailfill[30mm] 定高参——钉档 §三.3：无参在平衡栏呈「内容尾随框」，贴栏底须定高参；实测无参致末页平衡 Overfull\vbox 12.99pt，定高 30mm 三零） ----
\tailfill[30mm]

\end{multicols}
\end{document}
